\chapter{Geometryzacja rozmaitości trójwymiarowych}
\label{ch:geometrizacja-thurstona}
\index{geometryzacja Thurstona}
\index{geometrie Thurstona}

Rozmaitość trójwymiarowa może być złożona z obszarów o odmiennej
geometrii. Torus $T^3$ ma metrykę płaską, $S^3$ metrykę o dodatniej
krzywiźnie stałej, a istnieją zamknięte rozmaitości o krzywiźnie
stałej ujemnej. Twierdzenie geometryzacyjne mówi, w jakim sensie
takie modele wystarczają po rozcięciu dowolnej zamkniętej
rozmaitości trójwymiarowej. Poniżej podamy jego dokładną wersję topologiczną,
opiszemy osiem geometrii i wskażemy miejsce potoku Ricciego w dowodzie.
Głębokie wyniki przywołujemy ze źródeł; nie dowodzimy tutaj
twierdzenia geometryzacyjnego.

Przez \emph{zamkniętą} rozmaitość rozumiemy zwartą rozmaitość bez
brzegu. W całym rozdziale rozmaitości są gładkie i spójne, a w
sformułowaniu geometryzacji również orientowalne. Geometrie Kleina,
zdefiniowane przez jednorodne przestrzenie $G/H$, omówiliśmy szerzej
w Rozdziale~\ref{ch:geometria-kleina}. Tu potrzebny jest jeszcze
warunek metryczny: grupa symetrii ma zachowywać pewną metrykę
riemannowską.

\section{Od geometrii lokalnej do ilorazu}
\label{sec:thurston-model}

Geometria Kleina z Rozdziału~\ref{ch:geometria-kleina} opisuje przestrzeń
jednorodną przez działanie grupy: po wybraniu $x_0\in X$ mamy
$X\cong G/H$, gdzie $H=G_{x_0}$ jest stabilizatorem punktu.
Atlasy lokalnie modelowane na $X$, rozwinięcie i warunek zupełności
opisuje Sekcja~\ref{sec:thur-dop-struktury-gx}.
W wymiarze trzy interesują nas takie działania, z których można
otrzymać rozmaitości przez sklejenie punktów różniących się wybraną
symetrią. Warunki w następnej definicji odpowiadają kolejno za
metrykę, gładkość ilorazu i ograniczenie listy modeli.
Ponieważ grupy działają tu z lewej, ogólny iloraz orbit zapisujemy
$\Gamma\backslash X$; zapis $X/\Gamma$ w tradycyjnych nazwach
przykładów oznacza ten sam zbiór orbit.

\begin{definicja}[Geometria modelowa]
\label{def:thurston-model}
\index{geometria modelowa}
\index{geometria modelowa!maksymalna}
\emph{Geometrią modelową wymiaru trzy} nazywamy parę $(X,G)$,
w której $X$ jest jednospójną gładką rozmaitością wymiaru trzy,
a $G$ jest grupą Liego działającą na $X$ gładko, skutecznie
i przechodnio. Stabilizator $H=G_{x_0}$ pewnego, a więc każdego,
punktu jest zwarty. Ponadto istnieje dyskretna podgrupa
$\Gamma_0<G$, działająca na $X$ swobodnie, dla której
$\Gamma_0\backslash X$ jest zwarty. Mówimy wtedy, że model
\emph{dopuszcza zwarty iloraz}.

Geometria jest \emph{maksymalna}, gdy nie istnieje ściśle większa
grupa przekształceń $G'\subset\operatorname{Diff}(X)$ zawierająca
$G$, której działanie na tym samym $X$ spełnia powyższe warunki.
Porównujemy tu grupy jako grupy przekształceń $X$, a nie jedynie
ich abstrakcyjne izomorfizmy.
\end{definicja}

\subsection*{Stabilizator, zwartość i metryka}

Stabilizatory różnych punktów są sprzężone: jeśli $y=g\cdot x_0$,
to $G_y=gHg^{-1}$. Wystarczy więc sprawdzić zwartość jednego.
Zbiór jest zwarty, gdy z każdego jego pokrycia zbiorami otwartymi
można wybrać skończone podpokrycie. Na przykład
$\mathrm O(3)$ jest zwartym stabilizatorem zera w działaniu
$\R^3\rtimes\mathrm O(3)$ na $\R^3$, bo jest domknięty i ograniczony
w przestrzeni macierzy. Stabilizator $\mathrm{GL}(3,\R)$ w pełnej
grupie afinicznej nie jest zwarty: zawiera nieograniczony ciąg
$2^k I$. Działanie afiniczne nie zachowuje żadnej metryki
riemannowskiej, jak pokazuje Przykład~\ref{prz:klein-afiniczna}.

Zwarty stabilizator pozwala uśrednić dodatnio określony iloczyn
skalarny na $T_{x_0}X$ względem znormalizowanej miary Haara
i przenieść go przez $G$ do pozostałych punktów. Dokładny wzór,
niezależność od wyboru elementu przenoszącego i dowód zupełności
otrzymanej metryki są w
Stwierdzeniach~\ref{prop:klein-zwarty-stabilizator-metryka}
i~\ref{prop:klein-jednorodna-zupelna}.
Zwartość $H$ nie oznacza, że $X$ jest zwarta: w przykładzie
euklidesowym $H=\mathrm O(3)$ jest zwarta, lecz $X=\R^3$ nie jest.

\subsection*{Dyskretna grupa i gładki iloraz}

\begin{definicja}[Rozmaitość geometryczna]
\label{def:thurston-rozmaitosc-geometryczna}
\index{rozmaitość geometryczna}
\emph{Rozmaitością geometryczną} nazywamy gładką rozmaitość
$M=\Gamma\backslash X$ uzyskaną z maksymalnej geometrii modelowej $(X,G)$
przez iloraz względem dyskretnej podgrupy $\Gamma<G$ działającej
swobodnie. Iloraz jest \emph{zamkniętą rozmaitością geometryczną},
gdy jest ponadto zwarty.
\end{definicja}

Warunek istnienia $\Gamma_0$ w definicji modelu dotyczy
\emph{możliwości} otrzymania zwartego ilorazu; grupa $\Gamma$
wybrana dla konkretnej rozmaitości nie musi równać się $\Gamma_0$.
W $\R^3$ \emph{kratą pełnego rzędu} nazywamy podgrupę
$\Lambda=\Z v_1+\Z v_2+\Z v_3$ dla liniowo niezależnych
$v_1,v_2,v_3$; jej równoległościan podstawowy jest zwarty.
Na przykład dla modelu euklidesowego $\Gamma_0=\Z^3$ daje
zwarty torus $\R^3/\Z^3=T^3$, natomiast
$\Gamma=\Z(1,0,0)$ daje niezwarte $S^1\times\R^2$.
Kostka $[0,1]^3$ odwzorowuje się na cały $T^3$, co bezpośrednio
wyjaśnia jego zwartość. Samo $\Gamma=\{e\}$ jest dyskretne
i działa swobodnie, lecz ilorazem pozostaje niezwarte $\R^3$.

Związek warunków algebraicznych z geometrią jest następujący.
Dyskretność wyklucza grupy symetrii zbiegające do identyczności;
zwartość stabilizatora zapewnia właściwość działania $G$;
swoboda wyklucza punkty ustalone niejedynkowych elementów $\Gamma$.
Z tych trzech faktów wynika, że $\Gamma\backslash X$ jest gładką
rozmaitością, a rzut $X\to\Gamma\backslash X$ jest nakryciem i lokalną
izometrią dla wybranej metryki. Ponieważ $X$ jest jednospójna,
jest to nakrycie uniwersalne oraz
$\pi_1(\Gamma\backslash X)\cong\Gamma$; przy tym metryka ilorazowa
jest zupełna. Dowód twierdzenia o ilorazie znajduje się w
Twierdzeniu~\ref{prop:klein-dyskretny-iloraz}.
Jeśli grupa jest dyskretna, ale nie działa swobodnie,
iloraz może mieć punkty osobliwe: przekształcenie
$x\mapsto-x$ ustala $0\in\R^3$, a $\R^3/\{\pm I\}$
ma w obrazie zera otoczenie stożkowe o przekroju
$\mathbb{RP}^2$; rzut nie jest tam nakryciem. Jest to orbifold
z Przykładu~\ref{prz:klein-orbifoldy}, a jego ogólną definicję
podaje Definicja~\ref{def:klein-orbifold}.

Metryka na rozmaitości geometrycznej jest \emph{lokalnie
jednorodna}: każde dwa punkty mają małe otoczenia połączone
izometrią, bo można je podnieść do $X$ i przenieść działaniem
przechodnim $G$. Nie wymaga to, aby cała grupa izometrii ilorazu
działała na nim przechodnio. Przykładem zwartego ilorazu innego
niż torus jest $(S^2\times\R)/\Z\cong S^2\times S^1$,
gdzie $\Z$ przesuwa współrzędną rzeczywistą o liczby całkowite.
W geometrii sferycznej grupa $\{\pm1\}$ działa swobodnie
na $S^3$, a jej ilorazem jest $\mathbb{RP}^3$.
Wnętrze niektórych części powstałych po cięciu rozmaitości
ma natomiast niezwartą, zupełną metrykę hiperboliczną
o skończonej objętości; taki iloraz też jest rozmaitością
geometryczną w powyższym sensie, choć nie jest zamknięty.

Istnieje także droga w przeciwnym kierunku. Jeśli spójna
rozmaitość riemannowska ma zupełną metrykę lokalnie jednorodną,
to jej jednospójne nakrycie z podniesioną metryką jest jednorodne.
Jest to zewnętrzne twierdzenie o przedłużaniu lokalnych izometrii;
zob. \href{https://doi.org/10.1002/cpa.3160130408}
{I.\ M.\ Singer, \emph{Infinitesimally Homogeneous Spaces},
Comm. Pure Appl. Math. 13 (1960), 685--697} oraz omówienie
u \href{https://sschleimer.warwick.ac.uk/Teach/2021_3MFDS/References/1983the_geometries_of_three-manifolds.pdf}
{Scotta, wstęp i \S\,5}. Dla zamkniętej rozmaitości otrzymujemy
więc parę $(X,G)$ z $G=\operatorname{Isom}(X)$ i zwartym ilorazem.
Przekształcenia nakrycia $X\to M$ tworzą podgrupę
$\Gamma<G$ działającą swobodnie; jest ona dyskretna,
bo ciąg różnych przekształceń nakrycia zbieżny do identyczności
przenosiłby punkt $x$ dowolnie blisko niego, wbrew rozdzieleniu
arkuszy nad małym otoczeniem obrazu $x$.
Twierdzenie Myersa--Steenroda zapewnia, że $G$ jest grupą Liego
i działa właściwie, toteż stabilizator punktu jest zwarty
(\href{https://doi.org/10.2307/1968928}
{S.\ B.\ Myers i N.\ E.\ Steenrod,
\emph{The Group of Isometries of a Riemannian Manifold},
Ann. Math. 40 (1939), 400--416}). Warunek maksymalności
porządkuje dalej możliwe grupy przekształceń $X$.

\subsection*{Maksymalność modelu}

Maksymalność dotyczy grupy $G$ dla ustalonej przestrzeni $X$.
Grupa samych translacji $\R^3$ działa na $\R^3$ przechodnio
i ma zwarty stabilizator $\{e\}$ oraz zwarty iloraz przez $\Z^3$,
lecz jest zawarta w większej grupie
$E(3)=\R^3\rtimes\mathrm O(3)$, więc para $(\R^3,\R^3)$
nie jest maksymalna. Tak samo można dodać odbicia do
$\R^3\rtimes\mathrm{SO}(3)$.

\begin{przyklad}[Dlaczego pełny model euklidesowy jest maksymalny]
\label{prz:thurston-euklidesowa-maksymalna}
Załóżmy, że większa grupa $G'$ działająca na $\R^3$ zawiera
$E(3)$ i spełnia warunki Definicji~\ref{def:thurston-model}.
Zwarty stabilizator $G'_0$ daje $G'$-niezmienniczą metrykę
$g'$ przez uśrednianie. Ponieważ translacje należą do $E(3)$,
macierz $g'_x$ w zwykłych współrzędnych jest stała i równa
pewnej dodatnio określonej macierzy $Q$. Niezmienniczość
względem $\mathrm O(3)$ oznacza $A^{\mathsf T}QA=Q$ dla
każdego $A\in\mathrm O(3)$. Odbicia względem osi usuwają
wyrazy pozadiagonalne $Q$, a permutacje osi zrównują jego
wyrazy diagonalne. Zatem $Q=cI$ dla pewnego $c>0$.
Każdy element $G'$ zachowuje więc metrykę $c\,g_{\mathrm E}$.
Jej pełną grupą izometrii jest $E(3)$, jak wykazano w
Przykładzie~\ref{prz:klein-euklidesowa}. Stąd
$G'\subset E(3)$ i powiększenie jest niemożliwe.
\end{przyklad}

Nie wolno utożsamiać maksymalności z wyborem największej
podgrupy $\Gamma$ ani z jednoznacznością metryki. Jest to
warunek dotyczący \emph{grupy możliwych symetrii} na $X$.
Dwie pary $(X,G)$ i $(X',G')$ uznajemy za równoważne, gdy
istnieje dyfeomorfizm $f:X\to X'$ taki, że sprzężenie
$h\mapsto f\circ h\circ f^{-1}$ przenosi $G$ na $G'$.

\begin{twierdzenie}[Thurston: osiem maksymalnych geometrii]
\label{thm:thurston-osiem-modeli}
Każda maksymalna geometria modelowa wymiaru trzy jest równoważna
jednemu z ośmiu modeli opisanych w następnej sekcji.
\end{twierdzenie}
Jest to głęboki wynik zewnętrzny. Dokładne założenia i szkic
klasyfikacji podaje \href{https://sschleimer.warwick.ac.uk/Teach/2021_3MFDS/References/1983the_geometries_of_three-manifolds.pdf}
{P.\ Scott, \emph{The Geometries of 3-Manifolds},
Bull. London Math. Soc. 15 (1983), 401--487, \S\,5,
twierdzenie~5.1}; odrębność ośmiu typów wynika tam także
z twierdzenia~5.2. Torus $T^3$ i $\mathbb{RP}^3$ są ilorazami
modeli euklidesowego i sferycznego, a nie dodatkowymi modelami.

\section{Osiem geometrii}
\label{sec:thurston-eight}

Podajemy standardową metrykę lub wygodny jej przedstawiciel.
Zmiana wspólnej skali nie zmienia typu geometrii. Przez
$\operatorname{Isom}(X)$ oznaczamy grupę izometrii, a przez
$\operatorname{Isom}_0(X)$ jej składową jedności. Znaki krzywizn
poniżej dotyczą krzywizny sekcyjnej; w trzech ostatnich modelach
zależy ona od płaszczyzny stycznej. Opis i klasyfikacja są zgodne
z cytowaną pracą Scotta oraz z
\href{https://library.slmath.org/nonmsri/gt3m/PDF/Thurston-gt3m.pdf}
{W.\ P.\ Thurstonem, \emph{The Geometry and Topology of
Three-Manifolds}, zwłaszcza rozdz.\ 1--3}.

\subsection*{Trzy geometrie stałej krzywizny}

\textbf{Geometria sferyczna $S^3$.}
\index{geometrie Thurstona!sferyczna}
Sfera jednostkowa w $\R^4$ ma $K=1$ i pełną grupę izometrii
$\mathrm O(4)$, wymiaru $6$. Zamknięte ilorazy $S^3/\Gamma$,
gdzie skończona grupa $\Gamma<\mathrm O(4)$ działa swobodnie,
nazywają się \emph{sferycznymi formami przestrzennymi}.
Ich grupa podstawowa jest skończona.
Konkretnymi przykładami są przestrzenie soczewkowe
$L(p,q)$: dla $p\geq2$, $\gcd(p,q)=1$, grupa cykliczna rzędu $p$
działa na $S^3\subset\C^2$ przez
\[
 (z_1,z_2)\longmapsto
 (e^{2\pi i/p}z_1,e^{2\pi iq/p}z_2).
\]
Żadna nietrywialna potęga tego przekształcenia nie ma punktu
stałego: co najmniej jedna współrzędna jest niezerowa, a
$\gcd(p,q)=1$. Dla $p=2,q=1$ dostajemy $\mathbb{RP}^3$.

\textbf{Geometria euklidesowa $\mathbb E^3$.}
\index{geometrie Thurstona!euklidesowa}
Jest to $\R^3$ z $K=0$ i
$\operatorname{Isom}(\mathbb E^3)=\R^3\rtimes\mathrm O(3)$,
również wymiaru $6$. Dla kraty $\Lambda\cong\Z^3$ iloraz
$\R^3/\Lambda=T^3$ jest płaskim torusem. Twierdzenie Bieberbacha,
przywołane tu jako wynik zewnętrzny, mówi, że grupa podstawowa
każdej zamkniętej płaskiej rozmaitości trójwymiarowej zawiera $\Z^3$ jako
podgrupę skończonego indeksu. Źródło: cytowany Scott, \S\,4.

\textbf{Geometria hiperboliczna $\mathbb H^3$.}
\index{geometrie Thurstona!hiperboliczna}
W modelu górnej półprzestrzeni $z>0$ można wziąć metrykę
\[
 g=\frac{(\dd x)^2+(\dd y)^2+(\dd z)^2}{z^2},\qquad K=-1.
\]
Pełna grupa izometrii jest izomorficzna z grupą przekształceń
Lorentza zachowujących dodatnią powłokę hiperboloidy
w $\R^{3,1}$; jej wymiar to $6$. Istnieją zarówno zamknięte
ilorazy $\mathbb H^3/\Gamma$, jak i niezwarte ilorazy zupełne
o skończonej objętości. Do drugich należy dopełnienie węzła
ósemkowego w $S^3$; jego końce są \emph{ostrogami}, których
przekroje są torusami. Istnienie takiej metryki i twierdzenie
o hiperbolicznych wypełnieniach Dehna pochodzą od Thurstona
(\href{https://library.slmath.org/nonmsri/gt3m/PDF/Thurston-gt3m.pdf}
{\emph{The Geometry and Topology of Three-Manifolds}, rozdz.\ 4}).
Grupa podstawowa zamkniętego ilorazu jest nieskończona i nie jest
wirtualnie abelowa, w odróżnieniu od przypadku euklidesowego.
Konkretne sklejenia i działanie wypełnienia Dehna na grupie
podstawowej omówimy w Sekcji~\ref{sec:thur-dop-hiperboliczne}.

\subsection*{Dwa iloczyny powierzchni z prostą}

\textbf{Geometria $S^2\times\R$.}
\index{geometrie Thurstona!$S^2\times\R$}
Metryka iloczynowa ma krzywiznę $1$ na płaszczyznach stycznych
do $S^2$ i $0$ na płaszczyznach zawierających kierunek $\R$.
Grupą izometrii jest
$\operatorname{Isom}(S^2)\times\operatorname{Isom}(\R)$,
o składowej jedności wymiaru $4$.
Iloraz $S^2\times(\R/\Z)=S^2\times S^1$ jest zamknięty i ma
grupę podstawową $\Z$. Ten przykład jest \emph{pierwszy}, lecz
nie \emph{nierozkładalny}: sfera $S^2\times\{t\}$ nie ogranicza
kuli. Innym orientowalnym ilorazem tej geometrii jest
$\mathbb{RP}^3\#\mathbb{RP}^3$. Ta rozmaitość nie jest pierwsza:
jej dwa składniki pierwsze mają geometrię sferyczną.

\textbf{Geometria $\mathbb H^2\times\R$.}
\index{geometrie Thurstona!$\mathbb H^2\times\R$}
Metryka iloczynowa ma krzywiznę $-1$ na płaszczyznach stycznych
do $\mathbb H^2$ i $0$ na mieszanych. Grupa izometrii jest
$\operatorname{Isom}(\mathbb H^2)\times\operatorname{Isom}(\R)$,
ze składową jedności wymiaru $4$.
Jeżeli $\Sigma_g$ jest zamkniętą orientowalną powierzchnią
rodzaju $g\geq2$ z metryką hiperboliczną, to
$\Sigma_g\times S^1$ jest zamkniętym ilorazem tej geometrii.
Jego grupa podstawowa to $\pi_1(\Sigma_g)\times\Z$;
drugi czynnik opisuje regularne włókno okręgowe.
Istnienie metryki na każdej zamkniętej orientowalnej
$\Sigma_g$ o $g\geq2$ jest częścią twierdzenia o uniformizacji,
tu użytego jako wynik zewnętrzny; zob.
\href{https://library.slmath.org/nonmsri/gt3m/PDF/Thurston-gt3m.pdf}
{W.\ P.\ Thurston, \emph{The Geometry and Topology of
Three-Manifolds}, rozdz.\ 2}.

\subsection*{Trzy geometrie związane z grupami Liego}

\textbf{Geometria $\widetilde{\mathrm{SL}}(2,\R)$.}
\index{geometrie Thurstona!$\widetilde{\mathrm{SL}}(2,\R)$}
Przestrzeń stycznych wektorów jednostkowych do
$\mathbb H^2$ jest izomorficzna z $\mathrm{PSL}(2,\R)$;
jej nakrycie uniwersalne daje model
$\widetilde{\mathrm{SL}}(2,\R)$.
Podnosimy nań metrykę Sasakiego. Naturalne odwzorowanie
na $\mathbb H^2$ ma jednowymiarowe włókna, a składowa
jedności grupy izometrii ma wymiar $4$.
Krzywizna sekcyjna przyjmuje znaki dodatnie i ujemne;
nie jest to przestrzeń stałej krzywizny.
Dla hiperbolicznej powierzchni $\Sigma_g$ jej wiązka
jednostkowych wektorów stycznych $U\Sigma_g$ jest zamkniętym
ilorazem tego modelu. Jako wiązka okręgowa nad $\Sigma_g$
ma niezerową liczbę Eulera $\chi(\Sigma_g)=2-2g$.
To właśnie różni ją od iloczynu $\Sigma_g\times S^1$.
Jawną metrykę wyprowadzimy poniżej, a znak i wartość liczby
Eulera w Przykładzie~\ref{prz:thur-dop-seifert-euler}.

\begin{przyklad}[Jawna metryka geometrii $\widetilde{\mathrm{SL}}(2,\R)$]
\label{prz:thur-dop-tilde-sl-metryka}
Weźmy model półpłaszczyzny hiperbolicznej
$\mathbb H^2=\{(x,y):y>0\}$ z metryką
$h=y^{-2}(\dd x^2+\dd y^2)$ oraz zorientowaną
ramą ortonormalną $e_1=y\partial_x$, $e_2=y\partial_y$.
Wzory na symbole Christoffela z
Przykładu~\ref{prz:christoffel-hiperboliczna} dają
$\nabla e_1=A\,e_2$ i $\nabla e_2=-A\,e_1$,
gdzie $A=\dd x/y$ jest formą koneksji tej ramy.
Niech $\varphi$ oznacza kąt jednostkowego wektora
$v=\cos\varphi\,e_1+\sin\varphi\,e_2$.
W tych współrzędnych metryka Sasakiego na $U\mathbb H^2$
mierzy osobno prędkość punktu bazy i zmianę wektora
względem transportu równoległego. Wzdłuż krzywej $(p(s),v(s))$
mamy
$\nabla_{\dot p}v=(\dot\varphi+A(\dot p))
(-\sin\varphi\,e_1+\cos\varphi\,e_2)$.
Kwadrat normy prędkości to więc
$h(\dot p,\dot p)+(\dot\varphi+A(\dot p))^2$,
co daje
\begin{equation}\label{eq:thur-dop-tilde-sl-metryka}
 g=\frac{\dd x^2+\dd y^2}{y^2}
    +\left(\dd\varphi+\frac{\dd x}{y}\right)^2.
\end{equation}
Na wiązce jednostkowych wektorów kąt jest okresowy
($\varphi\in\R/2\pi\Z$); po przejściu do nakrycia
uniwersalnego rozwijamy go do $\varphi\in\R$,
uzyskując jedną z metryk modelu
$\widetilde{\mathrm{SL}}(2,\R)$.

Forma $\alpha=\dd\varphi+\dd x/y$ znika na
wektorach poziomych, a jej krzywizna wynosi
\[
 \dd\alpha=\frac{\dd x\wedge\dd y}{y^2}
     =\mathrm{vol}_{\mathbb H^2}.
\]
Jeżeli krzywa w bazie jest brzegiem obszaru $D$,
to poziome podniesienie zmienia kąt o
$\Delta\varphi=-\int_D\mathrm{vol}_{\mathbb H^2}$
(modulo $2\pi$ w wiązce $U\mathbb H^2$).
Istotnie, poziomość oznacza $\dot\varphi=-\dot x/y$;
całkujemy po brzegu i stosujemy Stokesa do $A$.
W iloczynie $\mathbb H^2\times\R$ mamy natomiast
płaską koneksję produktową $\dd\varphi$.
Właśnie niezerowa krzywizna formy $\alpha$ odróżnia
te dwie geometrie jednorodne. Dalszą konsekwencję
topologiczną obliczymy w
Przykładzie~\ref{prz:thur-dop-seifert-euler}.
\end{przyklad}

\textbf{Geometria $\mathrm{Nil}$.}
\index{geometrie Thurstona!Nil}
Wygodny model to rzeczywista grupa Heisenberga
\[
 N=\left\{\begin{pmatrix}1&x&z\\0&1&y\\0&0&1\end{pmatrix}
 :x,y,z\in\R\right\},
 \qquad
 g=(\dd x)^2+(\dd y)^2+(\dd z-x\,\dd y)^2.
\]
Metryka jest niezmiennicza względem lewych translacji.
W ortonormalnej bazie lewostronnie niezmienniczej $X,Y,Z$
jedyny niezerowy nawias bazowy to $[X,Y]=Z$. Wzór na koneksję
metryczną daje $K(X,Y)=-3/4$ oraz
$K(X,Z)=K(Y,Z)=1/4$. Pełna składowa jedności grupy
izometrii ma wymiar $4$: do translacji dochodzą obroty
płaszczyzny rozpiętej przez $X,Y$.
Macierze o całkowitych $x,y,z$ tworzą kratę $N_{\Z}$,
bo przy mnożeniu nowa współrzędna $z$ ma postać $z+z'+xy'$.
Iloraz lewostronny $N_{\Z}\backslash N$ jest zwartą wiązką
okręgową nad $T^2$
o niezerowej liczbie Eulera. Jego grupa podstawowa
$N_{\Z}$ jest nilpotentna, lecz nieabelowa.
Szczegółowy rachunek koneksji i krzywizn zawierają
Przykłady~\ref{prz:thur-dop-koneksja-nil}
i~\ref{prz:thur-dop-krzywizna-nil}.
Zwartość ilorazu można sprawdzić bez dodatkowego twierdzenia
o kratach. Dla $(x,y,z)$ wybieramy całkowite $m,n$ tak,
aby $x-m,y-n\in[0,1)$, a następnie całkowite $k$ tak,
aby $z-my+k\in[0,1)$. Lewa translacja przez $(-m,-n,k)$
przenosi punkt do kostki $[0,1]^3$. Jej obraz pokrywa więc
cały iloraz, który jest zwarty. Działanie jest swobodne,
gdyż lewa translacja ustalająca punkt musi być identycznością;
dyskretność wynika z całkowitości współrzędnych.

\begin{przyklad}[Rozkład kontaktowy i holonomia w geometrii Nil]
\label{prz:thur-dop-nil-kontakt}
Forma $\theta=\dd z-x\,\dd y$ występująca w metryce Nil
wyznacza rozkład dwuwymiarowych podprzestrzeni
$\mathcal H=\ker\theta\subset TN$.
Ponieważ $\dd\theta=-\dd x\wedge\dd y$, mamy
\[
 \theta\wedge\dd\theta=-\dd x\wedge\dd y\wedge\dd z\ne0.
\]
Rozkład jest zatem \emph{kontaktowy}: nie jest styczny
do żadnej lokalnej rodziny powierzchni. Istotnie, gdyby
$\mathcal H$ było styczne do powierzchni, to nawiasy pól
stycznych również byłyby styczne. Tymczasem dla
$E_1=\partial_x$ oraz $E_2=\partial_y+x\partial_z$
mamy $E_1,E_2\in\mathcal H$ i
$[E_1,E_2]=\partial_z\notin\mathcal H$.

Na $N\to\R^2$, $(x,y,z)\mapsto(x,y)$,
forma $\theta$ jest także formą koneksji poziomej
przy pionowych translacjach współrzędnej $z$.
Poziome podniesienie krzywej $\gamma(s)=(x(s),y(s))$
spełnia $\theta(\dot{\widetilde\gamma})=0$, czyli
$\dot z=x\dot y$. Dla poziomego podniesienia zamkniętej
krzywej podstawy przyrost współrzędnej pionowej wynosi
\begin{equation}\label{eq:thur-dop-nil-holonomia-pola}
 \Delta z=\oint_\gamma x\,\dd y
 =\int_D\dd x\wedge\dd y,
\end{equation}
jeżeli $\gamma=\partial D$ z orientacją brzegową.
Ostatnia równość jest twierdzeniem Stokesa.
W szczególności obieg dodatnio zorientowanego
kwadratu jednostkowego zmienia $z$ o $1$.
Po przejściu do ilorazu przez kratę całkowitą otrzymujemy
wiązkę okręgową nad $T^2$ o niezerowym skręceniu;
wartość bezwzględna jego liczby Eulera wynosi $1$.
Tak otrzymujemy geometryczny sens nieabelowego
komutatora w całkowitej grupie Heisenberga.
W ilorazie współrzędna włókna należy do $\R/\Z$:
przyrost $1$ oznacza tam powrót do tego samego punktu.
Jego całkowita wartość na podniesieniu rzeczywistym
zachowuje informację o skręceniu. Z generatorem okręgu
zapisanym jako kąt $2\pi z$ forma koneksji to $2\pi\theta$;
konwencja Eulera z Przykładu~\ref{prz:thur-dop-seifert-euler}
daje reprezentanta $-\dd\theta=\dd x\wedge\dd y$ na $T^2$,
o całce $1$.
\end{przyklad}

\begin{figure}[!htbp]
\centering
\begin{tikzpicture}[>=Stealth,line cap=round,line join=round]
 \fill[manifoldblue!12] (0,0) rectangle (3.0,2.0);
 \draw[->,manifoldblue!85!black,very thick] (0,0)--(3,0)
  node[midway,below,font=\small] {$\Delta z=0$};
 \draw[->,manifoldred,very thick] (3,0)--(3,2)
  node[midway,right,font=\small] {$\Delta z=ab$};
 \draw[->,manifoldblue!85!black,very thick] (3,2)--(0,2)
  node[midway,above,font=\small] {$\Delta z=0$};
 \draw[->,manifoldblue!85!black,very thick] (0,2)--(0,0)
  node[midway,left,font=\small] {$\Delta z=0$};
 \node[below left,font=\small] at (0,0) {$(0,0)$};
 \node[below right,font=\small] at (3,0) {$(a,0)$};
 \node[above right,font=\small] at (3,2) {$(a,b)$};
 \node[above left,font=\small] at (0,2) {$(0,b)$};
 \draw[manifoldblue!45!black] (6.3,-.25)--(6.3,2.4);
 \fill[manifoldblue!85!black] (6.3,.15) circle (2.2pt)
   node[right,font=\small] {$z_0$};
 \fill[manifoldred] (6.3,1.85) circle (2.2pt)
   node[right,font=\small] {$z_0+ab$};
 \draw[->,manifoldred,very thick] (6.05,.15)--(6.05,1.85)
   node[midway,left,font=\small] {$ab$};
 \node[below,font=\small] at (6.3,-.25) {włókno nad $(0,0)$};
\end{tikzpicture}
\caption{Poziome podniesienie dodatniego obiegu prostokąta
w Nil spełnia $\dd z=x\,\dd y$. Tylko bok $x=a$ wnosi $ab$.
Pętla w bazie zamyka się, ale jej podniesienie w $N$ zmienia wysokość.
Na rysunku $a,b>0$; we włóknie okręgowym przyrost liczymy modulo $1$.}
\label{fig:thur-dop-nil-holonomia}
\end{figure}

\textbf{Geometria $\mathrm{Sol}$.}
\index{geometrie Thurstona!Sol}
Niech $\R^2\rtimes\R$ ma mnożenie
\[
 (x,y,t)(x',y',t')=
 (x+e^t x',\;y+e^{-t}y',\;t+t')
\]
i lewostronnie niezmienniczą metrykę
\[
 g=e^{-2t}(\dd x)^2+e^{2t}(\dd y)^2+(\dd t)^2.
\]
Jedna pozioma oś rozszerza się przy wzroście $t$,
druga kurczy. Dla ortonormalnych kierunków $X,Y,T$
odpowiednio wzdłuż $x,y,t$ mamy
$K(X,Y)=1$ i $K(X,T)=K(Y,T)=-1$.
Rachunek tych krzywizn zawiera
Przykład~\ref{prz:thur-dop-krzywizna-sol}, a pierwsze całki
geodezyjnych Przykład~\ref{prz:thur-dop-geodezyjne-sol}.
Składowa jedności grupy izometrii ma wymiar $3$.
Przykład zwartego ilorazu otrzymujemy z macierzy
\[
 A=\begin{pmatrix}2&1\\1&1\end{pmatrix}\in\mathrm{SL}(2,\Z).
\]
Jej wartości własne są dodatnie, wzajemnie odwrotne i różne od $1$.
Torusowa wiązka nad $S^1$ z monodromią $A$ ma grupę podstawową
$\Z^2\rtimes_A\Z$ i geometrię $\mathrm{Sol}$: w bazie
wektorów własnych $A$ jej działanie na $\R^2$ odpowiada
przesunięciu współrzędnej $t$ o $\log\lambda$, gdzie
$\lambda>1$ jest większą wartością własną.
Grupa ta jest rozwiązalna, ale nie jest wirtualnie nilpotentna.

Sprawdźmy, że konstrukcja rzeczywiście daje zwarty iloraz
z podaną metryką. Wybierzmy macierz $P$ z wektorami własnymi
$A$ w kolumnach, tak aby
$D=P^{-1}AP=\operatorname{diag}(\lambda,\lambda^{-1})$,
gdzie $\lambda=(3+\sqrt5)/2$. Współrzędne $(x,y)$ są
współrzędnymi w tej bazie, więc krata ma postać
$\Lambda=P^{-1}\Z^2$, niekoniecznie $\Z^2$.
Mamy $D\Lambda=\Lambda$. Translacje o $\Lambda$ oraz
$F(x,y,t)=(\lambda x,\lambda^{-1}y,t+c)$,
gdzie $c=\log\lambda$, są izometriami. Dla $F$ wynika to z
$e^{-2(t+c)}\lambda^2=e^{-2t}$ i
$e^{2(t+c)}\lambda^{-2}=e^{2t}$.
Ponadto $F\circ T_\ell\circ F^{-1}=T_{D\ell}$,
więc generują grupę $\Lambda\rtimes_D\Z\cong\Z^2\rtimes_A\Z$.

Dla każdego elementu z niezerową potęgą $F$ współrzędna $t$
zmienia się o niezerową wielokrotność $c$, a pozostałe elementy
są translacjami kraty; działanie jest zatem swobodne.
Jest też właściwie nieciągłe: dla ustalonego zwartego zbioru
tylko skończenie wiele przesunięć o $kc$ może go spotkać,
a dla każdego takiego $k$ tylko skończenie wiele translacji kraty.
Równoległobok fundamentalny $\Lambda$ pomnożony przez $[0,c]$
odwzorowuje się na cały iloraz, co dowodzi zwartości.
Identyfikacja brzegów ma postać $(v,0)\sim(Dv,c)$.
Jeżeli monodromię zapisujemy zamiast tego regułą
$(v,c)\sim(f(v),0)$, jest nią $D^{-1}$, czyli $A^{-1}$
w bazie całkowitej. Odwrócenie kierunku okręgu bazy zamienia
$A$ z $A^{-1}$ i wyjaśnia obie konwencje torusa odwzorowania.

\begin{uwaga}[Porównanie ośmiu modeli]
\label{uw:thurston-eight-comparison}
Trzy przestrzenie $S^3,\mathbb E^3,\mathbb H^3$ mają stałą
krzywiznę i sześciowymiarową grupę izometrii. Dwa iloczyny
mają płaszczyzny o krzywiźnie $0$ oraz odpowiednio $+1$ i $-1$.
Modele $\widetilde{\mathrm{SL}}(2,\R)$, $\mathrm{Nil}$
i $\mathrm{Sol}$ mają oba znaki krzywizny. Ich typowe ilorazy
odróżnia topologia: niezerowe skręcenie wiązki okręgowej
nad powierzchnią hiperboliczną, niezerowe skręcenie nad torusem
oraz hiperboliczna monodromia wiązki torusowej.
Sama informacja, że $K$ zmienia znak, nie wystarcza więc
do odróżnienia trzech ostatnich geometrii.
\end{uwaga}

\section{Rozwłóknienia Seiferta i geometria}
\label{sec:thurston-seifert}
\index{rozwłóknienie Seiferta}

Orbifold pozwala zapisać dane punktów, nad którymi włókna
rozwłóknienia Seiferta zachowują się inaczej niż nad pozostałymi
punktami bazy. Potrzebujemy tutaj orientowalnych orbifoldów
powierzchniowych bez krawędzi zwierciadlanych.

\begin{definicja}[Orbifold powierzchniowy z punktami stożkowymi]
\label{def:thurston-orbifold-bazy}
\index{orbifold!powierzchniowy}
\index{punkt stożkowy}
Niech $\Sigma$ będzie zamkniętą orientowalną powierzchnią,
$b_1,\ldots,b_r\in\Sigma$ różnymi punktami, a $m_j\geq2$
liczbami całkowitymi. Orbifold $B$ ma przestrzeń podkładową
$|B|=\Sigma$, zwykłe mapy rozmaitości poza $b_j$ i mapę
$D/C_{m_j}$ przy $b_j$, gdzie $D\subset\C$ jest otwartym
dyskiem jednostkowym, a $C_{m_j}$ działa obrotami.
Punkt $b_j$ jest \emph{punktem stożkowym rzędu $m_j$}.
Piszemy $B=\Sigma(m_1,\ldots,m_r)$; dla $r=0$ jest to
zwykła powierzchnia traktowana jako orbifold.
\end{definicja}

Taki opis obejmuje wszystkie zamknięte orientowalne orbifoldy
powierzchniowe bez krawędzi zwierciadlanych: skończona podgrupa
$\mathrm{SO}(2)$ jest cykliczna, a jej niejedynkowe obroty
ustalają tylko środek dysku. Punkty stożkowe tworzą zbiór
domknięty i dyskretny, a taki podzbiór zwartej powierzchni
jest skończony.
W metryce otrzymanej z płaskiego dysku otoczenie $b_j$ ma
kąt stożkowy $2\pi/m_j$, choć topologicznie pozostaje dyskiem.
Zachowanie rzędu $m_j$ odróżnia $B$ od samej powierzchni $|B|$.

\begin{definicja}[Orbifoldowa charakterystyka Eulera]
\label{def:thurston-chi-orb}
\index{charakterystyka Eulera!orbifoldowa}
Wybierzmy triangulację $|B|$, w której wszystkie $b_j$ są
wierzchołkami, i oznaczmy przez $H_\sigma$ grupę lokalną
we wnętrzu komórki $\sigma$. Definiujemy
\[
 \chi^{\mathrm{orb}}(B)
 =\sum_{\sigma}\frac{(-1)^{\dim\sigma}}{|H_\sigma|}.
\]
Wszystkie komórki poza wierzchołkami $b_j$ mają wagę $1$,
a wierzchołek $b_j$ ma wagę $1/m_j$. Dlatego
\[
 \chi^{\mathrm{orb}}(B)
 =\chi(|B|)-\sum_{j=1}^r\left(1-\frac1{m_j}\right).
\]
Ostatni wzór pokazuje zarazem, że wynik nie zależy od
wybranej triangulacji.
\end{definicja}

\begin{przyklad}[Trzy znaki charakterystyki orbifoldowej]
\label{prz:thurston-chi-orb-znaki}
Dla $B=S^2(p,q,r)$ otrzymujemy bezpośrednio
\[
 \chi^{\mathrm{orb}}(B)=\frac1p+\frac1q+\frac1r-1.
\]
Stąd $S^2(2,3,5)$ ma charakterystykę $1/30$,
$S^2(2,3,6)$ ma charakterystykę $0$, a $S^2(2,3,7)$
ma charakterystykę $-1/42$. Znak tej liczby organizuje
tabelę geometrii Seiferta poniżej.

Inny przykład to $T^2/(x\mapsto-x)$. Inwolucja ma cztery
punkty stałe, których współrzędne na $\R^2/\Z^2$ należą do
$\{0,\tfrac12\}^2$. Kwadrat fundamentalny torusa po
identyfikacji przez półobrót tworzy \emph{poduszkę}, czyli
powierzchnię homeomorficzną ze sferą. Iloraz jest zatem
orbifoldem o przestrzeni podkładowej $S^2$ z czterema punktami
stożkowymi rzędu $2$. Zatem
$\chi^{\mathrm{orb}}\bigl(S^2(2,2,2,2)\bigr)
=2-4(1-1/2)=0$.
\end{przyklad}

\begin{definicja}[Dobry orbifold]
\label{def:thurston-orbifold-dobry}
\index{orbifold!dobry}
Orbifold jest \emph{dobry}, gdy ma nakrycie orbifoldowe przez
rozmaitość. W takim nakryciu lokalne odwzorowanie ma postać
$D^n/K\to D^n/H$ dla podgrupy $K\leq H$; gdy przestrzeń
nakrywająca jest rozmaitością, $K=\{e\}$. Rzut może być
rozgałęziony jako odwzorowanie przestrzeni podkładowych.
Orbifold bez takiego nakrycia nazywamy \emph{złym}.
Na przykład $S^2(2,2,2,2)$ jest dobry, bo otrzymaliśmy go
jako iloraz torusa, natomiast $S^2(p)$ dla $p\geq2$ jest zły.
\end{definicja}

Wyniki zewnętrzne Scotta mówią dokładniej, że wśród zamkniętych
orientowalnych orbifoldów powierzchniowych z punktami stożkowymi
jedyne złe przypadki to $S^2(p)$ oraz $S^2(p,q)$ dla $p\ne q$.
Każdy dobry zwarty orbifold powierzchniowy ma skończone
nakrycie przez powierzchnię i geometrię modelowaną na
$S^2$, $\mathbb E^2$ albo $\mathbb H^2$; w tym przypadku
znak $\chi^{\mathrm{orb}}$ odpowiada odpowiednio krzywiźnie
dodatniej, zerowej albo ujemnej. Są to twierdzenia
\href{https://sschleimer.warwick.ac.uk/Teach/2021_3MFDS/References/1983the_geometries_of_three-manifolds.pdf}
{P.\ Scotta, \emph{The Geometries of 3-Manifolds},
\S\,2, twierdzenia 2.3--2.5}; definicję lokalnych map i
ważonej charakterystyki podaje także
\href{https://library.slmath.org/nonmsri/gt3m/PDF/13.pdf}
{W.\ P.\ Thurston, \emph{The Geometry and Topology of
Three-Manifolds}, rozdział 13, \S\,13.2--13.3}.
Zła baza Seiferta może mimo to prowadzić do geometrycznej
rozmaitości całkowitej; tabela poniżej dotyczy rozmaitości,
a nie istnienia metryki stałej krzywizny na każdej wybranej bazie.

Rozmaitość Seiferta jest rozłożona na okręgi tak, że nad typowym
punktem powierzchni ilorazowej wygląda jak wiązka okręgowa.
Dokładniej, w przypadku orientowalnych włókien każde włókno ma
otoczenie będące ilorazem $D^2\times S^1$ przez działanie
cykliczne rzędu $m$ zadane wzorem
\[
 (z,u)\longmapsto
 (e^{2\pi i/m}z,e^{2\pi iq/m}u),\qquad \gcd(q,m)=1.
\]
Dla $m=1$ jest to zwykły iloczyn; dla $m>1$ włókno nad środkiem
dysku jest wyjątkowe, rzędu $m$. Działanie na przestrzeni
całkowitej jest swobodne: jeśli jego $k$-ta potęga ustala
$(z,u)$ dla $1\leq k<m$, to z $u\ne0$ wynika
$e^{2\pi iqk/m}=1$, a więc $m\mid qk$. Ponieważ
$\gcd(q,m)=1$, mielibyśmy $m\mid k$, sprzeczność.
Na dysku bazy ta sama potęga działa przez obrót i ustala
$z=0$. Stąd przestrzeń całkowita jest gładką rozmaitością,
a baza ma przy obrazie środka mapę orbifoldową $D/C_m$.
Poza wyjątkowymi włóknami lokalne grupy bazy są trywialne.
Liczba Eulera rozwłóknienia $e\in\Q$ mierzy jego skręcenie;
gdy nie ma włókien wyjątkowych, jest zwykłą całkowitą
liczbą Eulera wiązki okręgowej. Wystarcza nam rozróżnienie
$e=0$ i $e\ne0$, niezależne od konwencji znaku.
Na przykład wiązka iloczynowa $\Sigma_g\times S^1$ ma $e=0$,
natomiast $U\Sigma_g$ ma $e=\chi(\Sigma_g)\ne0$.

\begin{twierdzenie}[Geometria zamkniętej rozmaitości Seiferta]
\label{thm:thurston-seifert-table}
Załóżmy, że $M$ jest zamkniętą orientowalną rozmaitością
Seiferta z orientowalnym orbifoldem bazy $B$ oraz orientowalnymi
włóknami. Dla ustalonego rozwłóknienia typ geometrii $M$
wyznaczają znak $\chi^{\mathrm{orb}}(B)$ i warunek $e=0$:
\[
\begin{array}{c|cc}
 & e=0 & e\ne0\\ \hline
\chi^{\mathrm{orb}}(B)>0 & S^2\times\R & S^3\\
\chi^{\mathrm{orb}}(B)=0 & \mathbb E^3 & \mathrm{Nil}\\
\chi^{\mathrm{orb}}(B)<0 & \mathbb H^2\times\R
 & \widetilde{\mathrm{SL}}(2,\R)
\end{array}
\]
Jest to klasyfikacja dla podanych założeń, przywołana z
\href{https://sschleimer.warwick.ac.uk/Teach/2021_3MFDS/References/1983the_geometries_of_three-manifolds.pdf}
{P.\ Scotta, \emph{The Geometries of 3-Manifolds}, \S\,4,
tablica 4.1}. Przypadki z nieorientowalną bazą wymagają
odpowiednich wariantów i nie są objęte tą tabelą.
\end{twierdzenie}

Tabela wyjaśnia pary modeli o wspólnej bazie: nad powierzchnią
hiperboliczną brak skręcenia daje
$\mathbb H^2\times\R$, a skręcenie daje
$\widetilde{\mathrm{SL}}(2,\R)$. Nad torusem
otrzymujemy odpowiednio geometrię euklidesową i $\mathrm{Nil}$.
Geometrie hiperboliczna $\mathbb H^3$ i $\mathrm{Sol}$
nie występują w tej tabeli: ich zamknięte ilorazy nie są
rozmaitościami Seiferta. W szczególności wiązka torusowa
o hiperbolicznej monodromii z poprzedniej sekcji ma geometrię
$\mathrm{Sol}$, a nie $\mathrm{Nil}$.

\begin{przyklad}[Liczba Eulera wiązki jednostkowych wektorów]
\label{prz:thur-dop-seifert-euler}
Niech $\Sigma_g$ będzie zamkniętą, zorientowaną powierzchnią
rodzaju $g\ge2$ z metryką hiperboliczną. Rzut
$\pi:U\Sigma_g\to\Sigma_g$ jest główną wiązką
$\mathrm{SO}(2)$, ponieważ w każdym punkcie wektory
jednostkowe tworzą okrąg. Jej klasa Eulera jest klasą
Eulera zorientowanej wiązki stycznej $T\Sigma_g$.

W lokalnej zorientowanej ramie ortonormalnej $e_1,e_2$
zapiszmy $\nabla e_1=A\,e_2$ oraz
$\nabla e_2=-A\,e_1$. Z równania strukturalnego
koneksji na powierzchni otrzymujemy
$\dd A=-K\,\mathrm{vol}_{\Sigma_g}$
przy tej konwencji formy $A$; por.
\eqref{eq:gb-dalpha}. Przy obrocie ramy o lokalny kąt $\psi$
forma zmienia się na $A'=A+\dd\psi$, podczas gdy współrzędna
wektora zmienia się na $\varphi'=\varphi-\psi$.
Stąd $\dd\varphi'+A'=\dd\varphi+A$ jest globalną formą
koneksji na $U\Sigma_g$, a $\dd A$ skleja się do formy na bazie.

Zorientujmy włókno wzrostem kąta $\varphi$ od $e_1$ ku $e_2$.
Macierz koneksji w ramie $(e_1,e_2)$ to
$\bigl(\begin{smallmatrix}0&-A\\A&0\end{smallmatrix}\bigr)$,
więc składowa $\Omega_{12}$ macierzy krzywizny wynosi $-\dd A$.
Wzór Cherna--Weila na klasę Eulera z
Rozdziału~\ref{ch:chern-weil} daje reprezentanta
$-\dd A/(2\pi)$. Jest to klasa rzeczywista odpowiadająca
całkowitej klasie Eulera; do obliczenia jej liczby wystarcza całka.
Zatem
\[
 \langle e(T\Sigma_g),[\Sigma_g]\rangle
 =\frac1{2\pi}\int_{\Sigma_g}K\,\mathrm{vol}_{\Sigma_g}
 =\chi(\Sigma_g)=2-2g.
\]
Ostatnia równość jest twierdzeniem Gaussa--Bonneta
z Sekcji~\ref{sec:gauss-bonnet}, a nie osobnym postulatem
geometryzacji. W szczególności liczba Eulera
tej wiązki jest niezerowa, więc $U\Sigma_g$ nie jest
wiązką iloczynową $\Sigma_g\times S^1$.
To jest konkretny powód rozróżnienia w tabeli
\ref{thm:thurston-seifert-table} geometrii
$\widetilde{\mathrm{SL}}(2,\R)$ i $\mathbb H^2\times\R$.
\end{przyklad}

\section{Przykłady hiperboliczne: wielościany, ostrogi i wypełnienia}
\label{sec:thur-dop-hiperboliczne}
\index{wypełnienie Dehna}
\index{ostroga hiperboliczna}
Model $\mathbb H^3$ opisano już w sekcjach o metryce i o
geometrii Kleina. Teraz skupimy się na tym, jak z tego modelu
zbudować rozmaitości i jak obliczać ich niezmienniki.

\begin{definicja}[Wielościan fundamentalny i sklejanie ścian]
\label{def:thur-dop-wieloscian}
Niech $\Gamma<\operatorname{Isom}(\mathbb H^3)$ będzie grupą
dyskretną. Wielościan $P$ o niepustym wnętrzu nazywamy
\emph{wielościanem fundamentalnym}, gdy wnętrza $\gamma P$
są parami rozłączne dla różnych $\gamma\in\Gamma$,
a ich suma razem z brzegami pokrywa $\mathbb H^3$.
Zadanie odwrotne polega na zadaniu wielościanu i izometrii
parujących jego ściany, a następnie wykazaniu, że wyznaczają
one grupę dyskretną i właściwe sklejanie.
\end{definicja}

Nie wystarczy sam rysunek parowania ścian. W twierdzeniu
Poincarégo o wielościanach trzeba kontrolować między innymi
cykle krawędzi, sumy kątów i zgodność identyfikacji.
Wielościany \emph{idealne} mają wierzchołki na sferze
w nieskończoności $\partial_\infty\mathbb H^3$;
objętość może być skończona, mimo że wielościan
jest niezwarty w metryce hiperbolicznej.

\emph{Ostroga} jest końcem postaci $C\times[0,\infty)$,
gdzie $C$ jest zwartą płaską powierzchnią, a metryka ma postać
$\dd s^2+e^{-2s}h_C$. Ten wzór otrzymujemy z modelu
półprzestrzeni przez $z=e^s$ i iloraz poziomych horosfer.
Przekroje kurczą się jak $e^{-s}$, ale pionowa droga do
$s=\infty$ ma nieskończoną długość. Jednocześnie objętość
końca to $\operatorname{Area}(C,h_C)\int_0^\infty e^{-2s}\,\dd s
=\tfrac12\operatorname{Area}(C,h_C)$.
Wyjaśnia to, jak zupełny, niezwarty koniec może mieć skończoną objętość.

\begin{przyklad}[Czworościany idealne, Gieseking i węzeł ósemkowy]
\label{prz:thur-dop-gieseking-osemka}
Regularny czworościan idealny w $\mathbb H^3$ ma kąty
dwuścienne równe $\pi/3$ i skończoną objętość $v_3$.
Istnieje parowanie jego czterech ścian dające
\emph{rozmaitość Giesekinga}: jest ona nieorientowalna,
zupełna, niezwarta i ma skończoną objętość $v_3$.
Jej jedyna ostroga ma przekrój będący butelką Kleina.
Natomiast odpowiednie sklejenie \emph{dwóch} takich
czworościanów daje dopełnienie węzła ósemkowego
$M_8=S^3\setminus K_8$ o objętości $2v_3$.
Jest ono orientowalne, a jego jedyna ostroga ma przekrój torusowy.
Dopełnienie węzła ósemkowego jest orientowalnym
nakryciem dwukrotnym rozmaitości Giesekinga.

Z punktu widzenia geometrii nie wystarcza sprawdzenie,
że czworościany pokrywają rozmaitość topologiczną:
dookoła każdej klasy krawędzi suma kątów musi wynosić
$2\pi$, a identyfikacje wierzchołków idealnych
muszą dać zupełne przekroje euklidesowe.
Dla regularnego czworościanu sześć kątów po $\pi/3$
daje właśnie $2\pi$; w konkretnych sklejaniach należy
jeszcze sprawdzić wszystkie klasy krawędzi i wierzchołków.
Istnienie opisanych sklejeń i ich własności traktujemy
tutaj jako znane wyniki, nie jako konsekwencję samej sumy kątów.
Szczegóły parowań podaje
\href{https://www.mat.ufmg.br/posgrad/wp-content/uploads/TesesDissertacoes/Diss370.pdf}
{S.\ A.\ Pinillos Prado, \emph{Geometric structures on
3-dimensional manifolds}, \S\,3.1}; rozkład dopełnienia ósemki
omawia także Purcell, rozdziały 1 i 4.
\end{przyklad}


\begin{przyklad}[Zamknięta przestrzeń dwunastościenna Seiferta--Webera]
\label{prz:thur-dop-seifert-weber}
W poprzednim przykładzie otrzymaliśmy rozmaitości niezwarte
z wierzchołków idealnych. Istnieją także zwarte przykłady
skonstruowane z pojedynczego \emph{zwartego} wielościanu.
Weźmy foremny dwunastościan hiperboliczny o kącie dwuściennym
$2\pi/5$. Sklejamy każdą parę przeciwległych pięciokątnych
ścian, wykonując obrót o $3/10$ pełnego obrotu, czyli $3\pi/5$.
Pod tymi identyfikacjami trzydzieści krawędzi tworzy
sześć klas po pięć krawędzi. Suma kątów wokół każdej
krawędzi ilorazu wynosi więc $5(2\pi/5)=2\pi$.
Sklejenie jest zgodne również w wierzchołkach; jego
rezultatem jest zamknięta orientowalna rozmaitość
hiperboliczna, zwana przestrzenią Seiferta--Webera.
Kontrola otoczeń wierzchołków i pełnych warunków
sklejenia to istotna część konstrukcji, którą
przywołujemy tutaj jako znany wynik.
Nie należy mylić jej ze sferyczną przestrzenią
dwunastościenną Poincarégo: powstaje ona z innego
wielościanu i innej identyfikacji ścian.
Konstrukcja: cytowany Pinillos Prado, przykład 2.2.13.
\end{przyklad}

\begin{twierdzenie}[Sztywność Mostowa--Prasada]
\label{thm:thur-dop-mostow-prasad}
Niech $M$ i $N$ będą spójnymi zupełnymi rozmaitościami
hiperbolicznymi bez brzegu, o skończonej objętości i tym samym
wymiarze $n\ge3$. Każdy izomorfizm
$\pi_1(M)\cong\pi_1(N)$ jest realizowany przez izometrię,
z uwzględnieniem dróg potrzebnych do identyfikacji punktów bazowych.
W szczególności objętość hiperboliczna jest
niezmiennikiem topologicznym takich rozmaitości.
\end{twierdzenie}
Wynik ten cytujemy bez dowodu; zob.
\href{https://users.monash.edu/~jpurcell/book/HypKnotTheory.pdf}
{J.\ S.\ Purcell, \emph{Hyperbolic Knot Theory}, twierdzenie 6.1}.
Jego sens różni się od
sytuacji powierzchni hiperbolicznych, gdzie ustalona
topologia zazwyczaj dopuszcza nietrywialne deformacje
metryk o krzywiźnie stałej $-1$.

\begin{definicja}[Wypełnienie Dehna]
\label{def:thur-dop-dehn-fill}
Niech $N$ będzie spójną zwartą orientowalną rozmaitością wymiaru trzy
z brzegiem zawierającym torus $T$. Wybierzmy
uporządkowaną bazę $(\mu,\lambda)$ grupy
$H_1(T;\Z)\cong\Z^2$. \emph{Nachylenie} jest klasą
niezorientowanej krzywej prostej zamkniętej na $T$ o prymitywnej
klasie homologii $p\mu+q\lambda$, $\gcd(p,q)=1$.
Wypełnienie Dehna $N(p/q)$ otrzymujemy przez
przyklejenie pełnego torusa $D^2\times S^1$ do $T$
w taki sposób, aby obraz brzegu dysku meridianowego
$\partial D^2\times\{*\}$ przedstawiał wybrane nachylenie.
Pary $(p,q)$ i $(-p,-q)$ wyznaczają to samo nachylenie;
dla $q=0$ zapis $p/q$ oznacza $\infty$.
\end{definicja}

\begin{stwierdzenie}[Grupa podstawowa po wypełnieniu]
\label{prop:thur-dop-dehn-pi1}
Niech $N$ będzie jak wyżej,
a $\mu,\lambda$ oznaczają też elementy $\pi_1(N)$
wyznaczone przez ich obrazy w $N$, z ustalonymi
drogami bazowymi. Wówczas
\[
 \pi_1(N(p/q))\cong
 \pi_1(N)/\langle\!\langle\mu^p\lambda^q\rangle\!\rangle.
\]
\end{stwierdzenie}
\begin{proof}
Ponieważ krzywe $\mu,\lambda$ na torusie komutują,
klasa nachylenia ma w grupie postać $\mu^p\lambda^q$.
Wybierzmy $r,s\in\Z$ z $ps-qr=1$; istnieją dzięki
$\gcd(p,q)=1$. Krzywe $m=\mu^p\lambda^q$ i
$\ell=\mu^r\lambda^s$ tworzą bazę grupy torusa.
W dołączanym $D^2\times S^1$ krzywa $m$ staje się
ściągalnym meridianem, a $\ell$ przechodzi na generator $u$
grupy podstawowej pełnego torusa. Van Kampen daje więc
\[
 \pi_1(N(p/q))\cong
 (\pi_1(N)*\langle u\rangle)/
 \langle\!\langle m,\ell u^{-1}\rangle\!\rangle.
\]
Druga relacja pozwala wyeliminować $u=\ell$; zostaje
wyłącznie relacja $m=1$, czyli podany iloraz.
Zmiana dróg bazowych sprzęga element $m$, co nie zmienia
jego normalnego domknięcia. Nie zakładamy injektywności
$\pi_1(T)\to\pi_1(N)$.
\end{proof}

\begin{przyklad}[Homologia wypełnienia dopełnienia węzła]
\label{prz:thur-dop-dehn-h1}
Dla gładkiego węzła $K\subset S^3$ oznaczmy przez
$\nu K\cong D^2\times S^1$ jego domknięte otoczenie rurowe
i połóżmy $N=S^3\setminus\operatorname{int}\nu K$.
Mayer--Vietoris dla rozkładu $S^3=N\cup\nu K$
(po dodaniu małych otwartych kołnierzy) daje izomorfizm
$H_1(T^2;\Z)\cong H_1(N;\Z)\oplus H_1(\nu K;\Z)$,
bo $H_2(S^3;\Z)=H_1(S^3;\Z)=0$.
Meridian $\mu$ znika w pełnym torusie, a dowolna krzywa
$\lambda_0$ dopełniająca go do bazy przechodzi na generator
$H_1(\nu K;\Z)$. Izomorfizm pokazuje więc, że
$H_1(N;\Z)=\Z[\mu]$. Jeśli obraz $\lambda_0$ w tej grupie
wynosi $a[\mu]$, zastępujemy ją krzywą klasy
$\lambda=\lambda_0-a\mu$. Nadal tworzy ona bazę z $\mu$,
a w $H_1(N;\Z)$ znika. To \emph{preferowana długość}
(ang.\ \emph{longitude}). Po abelianizacji
relacji z poprzedniego stwierdzenia dostajemy
\[
 H_1(N(p/q);\Z)\cong
 \begin{cases}\Z/|p|\Z,&p\ne0,\\\Z,&p=0.\end{cases}
\]
Przykład pokazuje, jak operacja geometryczna zmienia
niezmiennik topologiczny. Nie wynika z niego jednak,
że dla $p=\pm1$ otrzymana rozmaitość jest sferą $S^3$:
zerowa pierwsza homologia nie rozstrzyga grupy podstawowej.
\end{przyklad}

\begin{twierdzenie}[Thurston: hiperboliczne wypełnienia Dehna]
\label{thm:thur-dop-dehn-hyperbolic}
Jeżeli $N$ jest spójną zwartą orientowalną rozmaitością wymiaru trzy,
$\partial N=T^2$, a jej wnętrze ma zupełną metrykę
hiperboliczną o skończonej objętości, to wszystkie poza
skończenie wieloma nachyleniami dają zamknięte
rozmaitości hiperboliczne $N(p/q)$.
\end{twierdzenie}
Twierdzenie to jest wynikiem zewnętrznym Thurstona,
a nie następstwem samej prezentacji grupy podstawowej.
Źródło: cytowana Purcell, twierdzenie 6.13 i \S\,6.4.
Z jednej niezwartej rozmaitości hiperbolicznej otrzymujemy
w ten sposób rodzinę zamkniętych przykładów.
Co więcej, przy nachyleniach, których długości na ustalonym
płaskim przekroju ostrogi dążą do nieskończoności, objętości
wypełnień dążą do objętości wyjściowej, pozostając od niej
ściśle mniejsze (Purcell, twierdzenia 6.29--6.30).
Takich objętości musi być nieskończenie wiele: skończony zbiór
liczb mniejszych od ustalonej liczby nie może się do niej zbiegać.
Sztywność Mostowa--Prasada daje zatem nieskończenie wiele
różnych typów topologicznych zamkniętych rozmaitości hiperbolicznych.

\begin{uwaga}[Związek z geometrią Sol]
Dla węzła ósemkowego nachylenie zerowe daje
szczególne, niehiperboliczne wypełnienie: torus odwzorowania
hiperbolicznego automorfizmu $T^2$, a więc rozmaitość
geometrii Sol. Włókno dopełnienia ósemki jest torusem
z usuniętym dyskiem. Wypełnienie według preferowanej długości
zamyka jego brzeg dyskiem i przedłuża monodromię do całego torusa;
na jego pierwszej homologii ma ona macierz
$\bigl(\begin{smallmatrix}2&1\\1&1\end{smallmatrix}\bigr)$,
z dokładnością do zmiany bazy i odwrócenia kierunku okręgu.
Korzystamy tu z opisu ósemki jako wiązki, a nie tylko z jej
pierwszej homologii; zob.
\href{https://library.slmath.org/nonmsri/gt3m/PDF/Thurston-gt3m.pdf}
{W.\ P.\ Thurston, \emph{The Geometry and Topology of
Three-Manifolds}, rozdział 4}.
To łączy geometrię węzłów z przykładem
monodromii macierzy hiperbolicznej w Sekcji~\ref{sec:thurston-eight}.
Sformułowanie o skończenie wielu wyjątkach nie pozwala jednak
samodzielnie ustalić listy wszystkich wyjątkowych nachyleń.
\end{uwaga}

\par\smallskip
{\footnotesize\noindent\emph{Dalsza lektura:}
\href{https://users.monash.edu/~jpurcell/book/HypKnotTheory.pdf}
{J.~S.~Purcell, \emph{Hyperbolic Knot Theory}};
\href{https://people.dm.unipi.it/martelli/Geometric_topology.pdf}
{B.~Martelli, \emph{An Introduction to Geometric Topology}},
rozdziały o geometriach Sol i hiperbolicznych rozmaitościach
z ostrogami;
\href{https://www.mat.ufmg.br/posgrad/wp-content/uploads/TesesDissertacoes/Diss370.pdf}
{S.~A.~Pinillos Prado, \emph{Geometric structures on
3-dimensional manifolds}} (2021), rozdziały 2--3.\par}

\section{Twierdzenie geometryzacyjne}
\label{sec:thurston-theorem}
\index{rozkład JSJ}
\index{rozkład pierwszy!rozmaitości trójwymiarowej}

Sfera potrzebna do \emph{sumy spójnej} i torus potrzebny do
\emph{cięcia JSJ} pełnią różne role. Zamknięta rozmaitość trójwymiarowa
jest \emph{pierwsza}, jeśli z zapisu $M=M_1\#M_2$ wynika,
że jeden składnik jest $S^3$. Jest \emph{nierozkładalna},
jeśli każda gładko zanurzona sfera $S^2$ ogranicza kulę $B^3$.
Zamknięte orientowalne składniki pierwsze są nierozkładalne
z wyjątkiem $S^2\times S^1$. Torus w nierozkładalnej
rozmaitości nazywamy \emph{nieściśliwym}, gdy inkluzja
$T^2\hookrightarrow M$ jest injektywna na $\pi_1$.
Równoważnie nie ma dysku ściskającego ten torus.
Jest to gładko zanurzony dysk $D\subset M$ taki, że
$D\cap T^2=\partial D$, a $\partial D$ nie ogranicza dysku
w $T^2$; południk torusa będącego brzegiem pełnego torusa
jest przykładem.
Kierunek od dysku do nieinjektywności jest bezpośredni;
kierunek odwrotny jest treścią zewnętrznego twierdzenia
o pętli Papakyriakopoulosa
(\href{https://doi.org/10.2307/1970113}
{\emph{On Dehn's Lemma and the Asphericity of Knots},
Ann. of Math. 66 (1957), 1--26}).
W tym sformułowaniu JSJ oznacza kanoniczne, po przyjęciu
zwykłych konwencji minimalności, cięcie wzdłuż rozłącznych
nieściśliwych torusów.

\begin{twierdzenie}[Geometryzacja Thurstona--Perelmana;
wynik zewnętrzny]
\label{thm:thurston-geometrization}
Niech $M$ będzie zamkniętą, spójną, orientowalną
rozmaitością trójwymiarową. Istnieje rozkład na sumę spójną skończenie
wielu rozmaitości pierwszych, jednoznaczny z pominięciem
składników $S^3$ i kolejności. Każdy składnik
nierozkładalny można rozciąć wzdłuż skończonej rodziny
rozłącznych nieściśliwych torusów JSJ tak, że każda
powstała część jest jednego z dwóch rodzajów:
\begin{enumerate}[label=\textup{(\roman*)}]
\item jej wnętrze ma zupełną metrykę hiperboliczną
  o stałej krzywiźnie $-1$ i skończonej objętości;
\item jest rozmaitością Seiferta.
\end{enumerate}
Składnik pierwszy $S^2\times S^1$ ma geometrię
$S^2\times\R$ i nie wymaga cięcia JSJ.
\end{twierdzenie}

Jest to wygodna wersja geometryzacji po połączeniu
twierdzenia o rozkładzie pierwszym, rozkładu JSJ
i dowodu hipotezy Thurstona. Pełne dowody i dokładniejsze
warianty, także z butelkami Kleina, podają
\href{https://www.claymath.org/wp-content/uploads/2022/03/Geom.pdf}
{J.\ Morgan i G.\ Tian, \emph{The Geometrization Conjecture}}
oraz \href{https://arxiv.org/pdf/0809.4040}
{ci sami autorzy, \emph{Completion of the Proof of the
Geometrization Conjecture}}. Trzeba odróżnić zdanie (i)
od twierdzenia, że każda część z brzegiem jest zwartym
ilorazem jednego z ośmiu modeli: jej \emph{wnętrze}
może być niezwartym ilorazem o skończonej objętości.
W przypadku (ii) część jest opisana przez rozwłóknienie
Seiferta; tabelę z poprzedniej sekcji stosuje się bezpośrednio
do zamkniętych rozmaitości spełniających jej założenia.
Rozmaitość \emph{geometryczna} jest pojedynczym ilorazem
$\Gamma\backslash X$ z Definicji~\ref{def:thurston-rozmaitosc-geometryczna}.
\emph{Geometryzacja} rozmaitości ogólnej oznacza opisanie jej
części po cięciach przez odpowiednie geometrie. Na przykład
$T^3\#T^3$ rozcina się wzdłuż sfery sumy spójnej na części,
które po dołączeniu kul dają dwa płaskie torusy; sama suma
spójna nie ma jednej globalnej geometrii modelowej
(zob. Scott, cytowana praca, wstęp).

\begin{przyklad}[Dwa proste wyniki cięcia]
\label{prz:thurston-cut}
Torus $T^3$ już sam jest zamkniętą rozmaitością geometryczną
typu euklidesowego;
nie trzeba go przecinać, mimo że zawiera wiele nieściśliwych
torusów. Dla rozmaitości sklejonej z dwóch części Seiferta
wzdłuż torusa włókna okręgowe mogą spotykać się pod różnymi
nachyleniami. Wtedy całość nie musi mieć jednej geometrii
modelowej, choć obie części są typu Seiferta. Rozkład JSJ
zachowuje informację o miejscu sklejenia.
\end{przyklad}

\section{Droga od potoku Ricciego do geometryzacji}
\label{sec:thurston-ricci-outline}

Rozdział~\ref{ch:potok-ricciego} wprowadza równanie
$\partial_tg=-2\operatorname{Ric}(g)$ i sprowadza istnienie
rozwiązania przez krótki czas do zewnętrznego twierdzenia
o układach parabolicznych. Aby uzyskać geometryzację,
potrzebne są cztery dalsze, głębokie kroki.

Po pierwsze, gdy krzywizna rośnie i gładki potok zbliża się
do osobliwości, analizuje się powiększone otoczenia punktów
o dużej krzywiźnie. Niezapadanie objętości i opis otoczeń
kanonicznych pozwalają rozpoznać cienkie szyjki oraz ich
zamknięte końce. Po drugie, chirurgia przecina wybrane szyjki,
zakleja końce i pozwala prowadzić potok dalej, jednocześnie
kontrolując zmianę topologii. Konstrukcję i jej analityczne
założenia opisuje
\href{https://arxiv.org/pdf/math/0211159}
{G.\ Perelman, \emph{The entropy formula for the Ricci flow and
its geometric applications}} oraz
\href{https://arxiv.org/pdf/math/0303109}
{G.\ Perelman, \emph{Ricci flow with surgery on three-manifolds}}.

Po trzecie, dla długiego czasu odpowiednio przeskalowany
potok rozdziela obszary \emph{grube}, z których powstają
hiperboliczne części o skończonej objętości, od obszarów
\emph{cienkich}, lokalnie zapadających się objętościowo.
Po czwarte, analiza tych ostatnich prowadzi do rozmaitości
grafowych, zbudowanych z części Seiferta. Ten ostatni krok
jest istotną częścią dowodu Morgana i Tiana, a nie
bezpośrednią konsekwencją samego równania potoku.

W przypadku jednospójnym istnieje krótsza droga do wniosku
topologicznego. Twierdzenie Perelmana o wygaśnięciu mówi:
jeśli zamknięta orientowalna rozmaitość trójwymiarowa ma rozkład
pierwszy bez składników asferycznych, to skonstruowany
potok z chirurgią wygasa w skończonym czasie dla każdej
gładkiej metryki początkowej
(\href{https://arxiv.org/pdf/math/0307245}
{G.\ Perelman, \emph{Finite extinction time for the solutions
to the Ricci flow on certain three-manifolds}}).
Dla rozmaitości jednospójnej warunek ten jest spełniony:
grupa podstawowa sumy spójnej jest iloczynem wolnym grup
składników, a zamknięta asferyczna rozmaitość trójwymiarowa nie jest
jednospójna. Gdyby była, byłaby ściągalna, czemu przeczy
jej niezerowa najwyższa homologia ze współczynnikami $\Z/2$.
Opis topologii usuwanych składników pokazuje
wtedy, że rozmaitość jest sumą spójną sferycznych form
przestrzennych i składników $S^2\times S^1$.
Jej grupa podstawowa jest iloczynem wolnym odpowiednich
grup skończonych i kopii $\Z$. Jeśli jest trywialna,
wszystkie te czynniki są trywialne; pozostaje $S^3$.
Ścisłe przejście przez chirurgię i klasyfikację usuwanych
składników znajduje się w
\href{https://www.claymath.org/wp-content/uploads/2022/03/Ricci-pdf.pdf}
{J.\ Morgana i G.\ Tiana, \emph{Ricci Flow and the Poincaré
Conjecture}}. Sam wniosek i jego miejsce wśród wersji
twierdzenia Poincarégo omawia Rozdział~\ref{ch:poincare-trzy}.

Powyższy łańcuch jest \emph{opisem idei dowodu}, a nie
dowodem twierdzenia geometryzacyjnego ani twierdzenia
Poincarégo. Analiza otoczeń osobliwych, chirurgii i zapadania
się wymaga osobnych twierdzeń, których założenia i dowody
znajdują się w podanych źródłach.

\par\smallskip
{\footnotesize\noindent\emph{Dalsza lektura:}
\href{https://library.slmath.org/nonmsri/gt3m/PDF/Thurston-gt3m.pdf}
{W.\ P.\ Thurston, \emph{The Geometry and Topology of
Three-Manifolds}}, geometrie modelowe i rozkłady;
\href{https://arxiv.org/pdf/0809.4040}
{J.\ Morgan i G.\ Tian, \emph{Completion of the Proof of the
Geometrization Conjecture}}, dokończenie dowodu geometryzacji.\par}
